\scanpage{24}
qui contient $\hat{\mathcal E}$ comme sous-extension. On désigne par $\mathord{\text{\begin{tikzpicture}[
  x=1em,y=1ex,baseline=0pt,
  line width=.045em,line cap=butt,line join=miter]
  \draw (.04,0)--(.04,1.60)--(.58,1.60)--(.58,1.40);
  \draw (.16,0)--(.16,1.40)--(.58,1.40);
\end{tikzpicture}}}$ le commutant de $\unsure{\mathfrak S_3|\mathcal S_3}$ dans $\widehat{\widehat\Gamma}$ \inserted{(\unsure{qui contient} $\dot\tau$)}, et par $\mathbb E$ l'extension correspondante de $\mathord{\text{\begin{tikzpicture}[
  x=1em,y=1ex,baseline=0pt,
  line width=.045em,line cap=butt,line join=miter]
  \draw (.04,0)--(.04,1.60)--(.58,1.60)--(.58,1.40);
  \draw (.16,0)--(.16,1.40)--(.58,1.40);
\end{tikzpicture}}}$ par $\hat\Pi$. Notons déjà que
\[
\left\{\begin{aligned}
\mathord{\text{\begin{tikzpicture}[
  x=1em,y=1ex,baseline=0pt,
  line width=.045em,line cap=butt,line join=miter]
  \draw (.04,0)--(.04,1.60)--(.58,1.60)--(.58,1.40);
  \draw (.16,0)--(.16,1.40)--(.58,1.40);
\end{tikzpicture}}}\cap\underset{\deleted{\unsure{\mathfrak S_3}}}{\unsure{\Gamma_{\unsure{0,3}}|\mathord{\text{\begin{tikzpicture}[
  x=1em,y=1ex,baseline=0pt,
  line width=.045em,line cap=butt,line join=miter]
  \draw (.04,0)--(.04,1.60)--(.58,1.60)--(.58,1.40);
  \draw (.16,0)--(.16,1.40)--(.58,1.40);
\end{tikzpicture}}}}}&=\{1,\tau\}\\
\mathord{\text{\begin{tikzpicture}[
  x=1em,y=1ex,baseline=0pt,
  line width=.045em,line cap=butt,line join=miter]
  \draw (.04,0)--(.04,1.60)--(.58,1.60)--(.58,1.40);
  \draw (.16,0)--(.16,1.40)--(.58,1.40);
\end{tikzpicture}}}\cap\unsure{\mathfrak S_3|\mathcal S_3}&=\{1\}
\end{aligned}\right.
\]
\margin{\deleted{\uline{\unsure{NB}}}}
\margin{\diagram{\begin{tikzcd}
&\unsure{\mathfrak S_3|\mathcal S_3}\times\widehat{\widehat\Gamma}^!\arrow[d,phantom,"{\cup}"]&&\\
\text{\uline{NB}}&\widehat{\widehat\Gamma}\arrow[r,phantom,"\supset"]&\underset{\deleted{\unsure{\Gamma}}}{\unsure{\mathfrak S_3|\mathcal S_3}}\times\mathord{\text{\begin{tikzpicture}[
  x=1em,y=1ex,baseline=0pt,
  line width=.045em,line cap=butt,line join=miter]
  \draw (.04,0)--(.04,1.60)--(.58,1.60)--(.58,1.40);
  \draw (.16,0)--(.16,1.40)--(.58,1.40);
\end{tikzpicture}}}&
\end{tikzcd}}
--- j'ignore si cette inclusion est une égalité, i.e. si $\mathord{\text{\begin{tikzpicture}[
  x=1em,y=1ex,baseline=0pt,
  line width=.045em,line cap=butt,line join=miter]
  \draw (.04,0)--(.04,1.60)--(.58,1.60)--(.58,1.40);
  \draw (.16,0)--(.16,1.40)--(.58,1.40);
\end{tikzpicture}}}\subset\widehat{\widehat\Gamma}^!$ est une égalité}
On veut donner d'abord une description plus symétrique des élts de $\widehat{\widehat{\mathcal E}}^!$, sans faire jouer un rôle particulier à $L_0$ engendré par $l_0$ $(L_0=\unsure{\hat k_0|\hat\kappa_0}(\hat{\mathbb Z}))$. Un $u\in\widehat{\widehat{\mathcal E}}$ est \inserted{pour être} défini par un multiplicateur
\[
(19)\qquad p=\chi(u)\in\hat{\mathbb Z}^{*}
\]
et des
\[
(20)\qquad g_i\in\hat\Pi\qquad i=0,1,\infty
\]
pour
\[
(21)\qquad u(l_i)=\operatorname{int}(g_i^{\unsure{\deleted{\unsure{-1}}|\unsure{-1}}})l_i^p
\]
les données (19), (20) étant liées par la relation
\[
(22)\qquad(\operatorname{int}(g_\infty^{\unsure{\deleted{\unsure{-1}}|\unsure{-1}}})l_\infty^p)
(\operatorname{int}g_1^{\unsure{\deleted{\unsure{-1}}|\unsure{-1}}}\cdot l_1^p)
(\operatorname{int}(g_0^{\unsure{\deleted{\unsure{-1}}|\unsure{-1}}})l_0^p)=1
\]
et par le fait que les $u(l_i)$ définis par (21) \uline{engendrent} $\hat\Pi$ \unsure{topolog\textsuperscript{t}|topologiquement}. Pour $u$ fixé, $p$
\mspage{23}
